\section{模型优缺点评价}

\subsection{模型优点}

本文围绕霍尔木兹海峡封锁对国际原油价格的影响，建立传统供需基准模型、短期冲击模型和中长期油价调节模型。模型的主要优点如下。

\begin{enumerate}
  \item 模型链条完整。传统供需基准给出低弹性条件下的反事实压力，短期冲击模型解释 0--60 天价格平台，中长期油价调节模型刻画 60--180 天情景路径和尾部风险。
  \item 变量口径清晰。附件价格只用于拟合和误差评价，EIA、JODI、OPEC、GPR 和 OVX 等外部数据只用于数量级约束和风险校验，不反向替代真实价格。
  \item 检验维度较全。短期模型与朴素上一日基准和滚动 ARIMA 基准比较，并通过滞后平移、机制消融、参数扰动和蒙特卡洛情景树检验模型稳定性。
  \item 结果解释性较强。模型没有只追求时间序列拟合，而是把供应缺口、SPR、商业库存、绕道运输、需求弹性和风险溢价放入统一的动态递推框架。
\end{enumerate}

\subsection{模型不足}

\begin{enumerate}
  \item 长期模型属于条件情景预测，不是对真实未来油价的精确点预测。若封锁持续期间出现新的军事、外交或金融冲击，状态转移概率仍需要重新校准。
  \item 外部数据存在频率和口径差异。EIA 为美国周度数据，JODI 为多国月度上报数据，OPEC 为正常情景供需基线，因此它们只能约束参数范围，不能完全替代全球日度真实供需数据。
  \item 风险变量仍有扩展空间。本文接入 GPR 和 OVX，但油轮保险费、BDTI/TD3C、期货期限结构和更完整的 IEA MODS 数据尚未形成可复现日度样本。
\end{enumerate}

\subsection{结论}

在赛题低短期弹性设定下，传统供需基准给出 278--337 美元/桶的压力上界，说明仅依赖静态供需关系会显著高估现实价格。加入 SPR、商业库存、绕道运输、需求收缩、恐慌衰减、地缘风险溢价和预期修复后，短期冲击模型能够解释冲突窗口内约 110--120 美元/桶的价格平台，并优于朴素上一日基准和滚动 ARIMA 基准。

中长期油价调节模型显示，中性情景第 180 天价格约为 97.82 美元/桶；悲观情景在供应中断偏大、SPR 等效预算耗尽、绕道恢复缓慢和制度风险溢价持续时，第 180 天价格约为 121.64 美元/桶，并保留二次跳涨风险。敏感性分析表明，封锁风险衰减速度、SPR 释放上限、供应中断量和绕道运输能力是长期结果的主要驱动因素。
